\section{附录：符号与术语表}\label{sec:appendix}

\begin{table}[htbp]
    \centering
    \caption{本文使用的主要符号}
    \label{tab:notation}
    \begin{tabular}{ll}
        \toprule
        符号 & 含义 \\
        \midrule
        $G$ & 约化群（如 $\operatorname{GL}_n$） \\
        $^LG$ & $G$ 的朗兰兹对偶群（$L$-群） \\
        $\mathbb{A}$ & $\mathbb{Q}$ 的 adèle 环；$G(\mathbb{A})$ 为 adelic 点 \\
        $\pi$ & 尖点自守表示；$\pi_v$ 为其局部因子 \\
        $\rho$ & 伽罗瓦表示 $G_K \to \operatorname{GL}_n$ \\
        $G_K$ & 绝对伽罗瓦群 $\operatorname{Gal}(\overline{K}/K)$ \\
        $W_F$, $W_F'$ & 局部域 $F$ 的韦伊群与 Weil--Deligne 群 \\
        $L(s,\pi)$, $L(s,\rho)$ & 自守 $L$ 函数与阿廷 $L$ 函数 \\
        $L(s,\pi,r)$ & 与对偶群表示 $r$ 配对的自守 $L$ 函数 \\
        $\operatorname{Sym}^k$ & $k$ 次对称幂表示（对称幂函子性） \\
        $R$, $T$ & 伽罗瓦形变环与 Hecke 代数（$R=T$） \\
        $X/\mathbb{F}_q$ & 有限域上的光滑射影曲线（函数域） \\
        $\operatorname{Bun}_G$ & 曲线上 $G$-主丛的叠（模空间） \\
        $\operatorname{Bun}_{G}(X)$ 上的 D-模 & 几何朗兰兹的自守一侧范畴 \\
        Shtuka & 带自相似映射的主丛（Drinfeld） \\
        \bottomrule
    \end{tabular}
\end{table}

\section*{术语对照}
\begin{spacing}{1.0}
\begin{itemize}[itemsep=0pt,topsep=4pt]
    \item 互反律（reciprocity）：伽罗瓦表示与自守表示按 $L$ 函数相等的对应；
    \item 函子性（functoriality）：$L$-同态诱导的自守表示之间的转移；
    \item 对偶群 / $L$-群（dual group / $L$-group $^LG$）：约化群的复对偶与其上的伽罗瓦作用之半直积；
    \item 尖点表示（cuspidal representation）：自守谱中“本原”的不可约成分；
    \item Satake 参数（Satake parameters）：非分歧局部表示的 $L$-打包数据；
    \item 内窥（endoscopy）：稳定轨道与普通轨道之偏差的记账群；
    \item 基本引理（fundamental lemma）：内窥转移的轨道积分恒等式（Ng\^o 定理）；
    \item 模性提升（modularity lifting）：由模 $p$ 自守性提升到 $\ell$-adic 自守性（$R=T$）；
    \item 势自守性（potential automorphy）：在有限个辅助域换基后自守的性质；
    \item Shtuka：函数域上带 Frobenius 型自相似结构的主丛，朗兰兹对应的几何载体；
    \item Hecke 特征单子（Hecke eigensheaf）：几何朗兰兹中尖点表示的范畴化对应物；
    \item 几何化（geometrization）：Fargues--Scholze 在 Fargues--Fontaine 曲线上重建局部朗兰兹的方案；
    \item 韦伊纲领 II（Weil II / arithmetization）：Scholze 提出的把几何朗兰兹算术化以证明数域朗兰兹的构想。
\end{itemize}
\end{spacing}
